\documentclass[11pt]{article}
\usepackage[margin=1in]{geometry}
\usepackage{amsmath}
\usepackage{booktabs}

\title{Exponential Feature Selection in the Data-Adaptive Mondrian Forest}
\date{\today}

\begin{document}
\maketitle

\section{Idea}

In Mondrian-CART the split feature is drawn with $p_j \propto s_j$, where $s_j$ is the expected CART
variance reduction of a uniformly random split along feature $j$. When $d$ is large, the many noise
features have small but nonzero scores, and together they take most of the probability, especially
in deep cells. We replace this rule by
\[
  p_j \;\propto\; \exp\!\Big(\beta\,\frac{s_j}{\max_k s_k}\Big), \qquad \beta = 10
  \qquad (\text{Mondrian-CART-exp}).
\]
The exponential widens the gap between high- and low-score features, while features with equal scores
still get equal probability. Scores are divided by the cell's largest score because raw scores shrink
with the number of points in the cell; without this, $e^{s_j}$ would be nearly uniform in deep cells.
The exponential clock is unchanged.

\paragraph{Setup.} $X \sim \mathrm{Unif}[0,1]^d$,
$y = \sin(4u_0) + 2(u_1 - 0.5)^2 + u_2 + \varepsilon$, with $\varepsilon \sim N(0, 0.1^2)$. In the
\emph{axis-aligned} setting, $u_0, u_1, u_2$ are the first three columns of $X$. In the \emph{rotated}
setting, they are three random orthonormal projections of $X$, so every column carries a little of the
signal. We use $n_{\text{train}} = 800$, 10 seeds and $M = 50$ trees. The lifetime of each Mondrian
forest is tuned on a validation set; for RF, \texttt{max\_features} and \texttt{min\_samples\_leaf} are
tuned on the same set. Mondrian-CART-null (CART score minus the noise baseline
$\mathrm{SSE}/(n-1)$, clipped at 0) is a separate fix of Mondrian-CART and is compared side by side.

\section{Results}

\paragraph{Where the splits go.} In the axis-aligned setting with $d = 50$ (chance level 6\%), the share
of splits on the 3 relevant features is 12\% for CART, 18\% for CART-null and 28\% for CART-exp
(5 seeds, 10 trees).

\paragraph{Test MSE.} Mean over 10 seeds, $M = 50$; the best Mondrian variant in each column is in bold.

\begin{center}
\begin{tabular}{lcccc}
\toprule
 & \multicolumn{2}{c}{Axis-aligned} & \multicolumn{2}{c}{Rotated} \\
\cmidrule(lr){2-3} \cmidrule(lr){4-5}
Method & $d=10$ & $d=50$ & $d=10$ & $d=50$ \\
\midrule
Mondrian-CART      & 0.0270          & 0.0588          & 0.1622          & 0.4035 \\
Mondrian-CART-null & \textbf{0.0232} & 0.0358          & 0.1477          & 0.3810 \\
Mondrian-CART-exp  & 0.0263          & \textbf{0.0310} & \textbf{0.1439} & \textbf{0.3712} \\
RF (tuned)         & 0.0259          & 0.0329          & 0.1696          & 0.3714 \\
\bottomrule
\end{tabular}
\end{center}

\section{Conclusions}

\begin{itemize}
  \item \textbf{CART-exp improves on CART} in every setting and every $d \in \{10, 20, 30, 50\}$.
  \item \textbf{CART-exp vs.\ CART-null.} In the rotated setting, CART-exp is better for every $d$.
    In the axis-aligned setting, CART-null is slightly better for $d \le 20$, the two tie at $d = 30$,
    and CART-exp is better at $d = 50$. The advantage of CART-exp grows with $d$.
  \item \textbf{CART-exp vs.\ RF.} CART-exp matches or beats RF in all cases. It ties at
    axis-aligned $d = 10$ and rotated $d = 50$, and is better everywhere else.
\end{itemize}
All differences stated above are at least 2 paired standard errors, except those described as ties.

\section{Caveats and next steps}

\begin{itemize}
  \item Only $\beta = 10$ has been tested. A smaller $\beta$ can be worse than proportional sampling:
    with $\beta = 3$, the minimum weight $e^{-\beta} \approx 5\%$ of the top feature is larger than
    what proportional sampling gives noise features in large cells.
  \item CART-exp often picks the largest lifetime in the grid (up to 6 of 10 seeds). We will extend
    the grid, or add $\lambda = \infty$, to check whether this limits its performance at small $d$.
\end{itemize}

\end{document}
